% --- Archon physics formalization source begin ---
% archon:physics
% archon:covers PhyXMiniProblems/problem_phyx_mini_0018.lean
% archon:source-report reports/phyx_mini/problem_phyx_mini_0018.source.json

\chapter{Physics problem phyx\_mini\_0018}
\label{ch:PhyXMiniProblems_problem_phyx_mini_0018}

\paragraph{Problem source.}
\#\# Physical scenario

A material having an index of refraction \textbackslash{}( n \textbackslash{}) is surrounded by vacuum and is in the shape of a quarter circle of radius \textbackslash{}( R \textbackslash{}) as shown in figure. A light ray parallel to the base of the material is incident from the left at a distance \textbackslash{}( L \textbackslash{}) above the base and emerges from the material at the angle \textbackslash{}( \textbackslash{}theta \textbackslash{}).

\#\# Question

Determine an expression for \textbackslash{}( \textbackslash{}theta \textbackslash{}) in terms of \textbackslash{}( n \textbackslash{}), \textbackslash{}( R \textbackslash{}), and \textbackslash{}( L \textbackslash{}).

\#\# Answer choices

- A: A: \textbackslash{}( \textbackslash{}sin\textasciicircum{}\{-1\}\textbackslash{}left[n\textbackslash{}sin\textbackslash{}left(\textbackslash{}tan\textasciicircum{}\{-1\}\textbackslash{}frac\{L\}\{R\} - \textbackslash{}sin\textasciicircum{}\{-1\}\textbackslash{}frac\{L\}\{nR\}\textbackslash{}right)\textbackslash{}right] \textbackslash{})

- B: B: \textbackslash{}( \textbackslash{}sin\textasciicircum{}\{-1\}\textbackslash{}left[n\textbackslash{}cos\textbackslash{}left(\textbackslash{}sin\textasciicircum{}\{-1\}\textbackslash{}frac\{L\}\{R\} - \textbackslash{}sin\textasciicircum{}\{-1\}\textbackslash{}frac\{L\}\{nR\}\textbackslash{}right)\textbackslash{}right] \textbackslash{})

- C: C: \textbackslash{}( \textbackslash{}sin\textasciicircum{}\{-1\}\textbackslash{}left[n\textbackslash{}sin\textbackslash{}left(\textbackslash{}sin\textasciicircum{}\{-1\}\textbackslash{}frac\{L\}\{R\} - \textbackslash{}sin\textasciicircum{}\{-1\}\textbackslash{}frac\{L\}\{nR\}\textbackslash{}right)\textbackslash{}right] \textbackslash{})

- D: D: \textbackslash{}( \textbackslash{}sin\textasciicircum{}\{-1\}\textbackslash{}left[n\textbackslash{}sin\textbackslash{}left(\textbackslash{}sin\textasciicircum{}\{-1\}\textbackslash{}frac\{L\}\{R\} - \textbackslash{}cos\textasciicircum{}\{-1\}\textbackslash{}frac\{L\}\{nR\}\textbackslash{}right)\textbackslash{}right] \textbackslash{})

\#\# Figure caption (auxiliary; use the image as primary evidence)

The image illustrates the refraction of light through a semi-circular lens or prism.

1. **Shapes and Objects:**
   - A semi-circular lens is depicted, with the straight edge on the right side and the curved edge on the left.
   - The lens is labeled with the refractive index "n."

2. **Rays:**
   - An "Incoming ray" enters the lens from the left, parallel to a horizontal base line.
   - The ray refracts at the curved surface and changes direction inside the lens.
   - The "Outgoing ray" exits the lens on the right side, diverging downwards.

3. **Labels and Measurements:**
   - The vertical distance from the base line to the incoming ray is labeled "L."
   - The radius of the semi-circle is labeled "R."
   - The angle of refraction inside the lens is labeled with the Greek letter "θ."

4. **Lines:**
   - Dashed lines are used to indicate the normal at the point of refraction and to extend the angle "θ."

The diagram is likely used to explain the principles of refraction and the behavior of light as it passes through a medium with a different refractive index.

\#\# Dataset metadata

Geometrical Optics; Physical Model Grounding Reasoning, Spatial Relation Reasoning

\paragraph{Recorded answer/context.}
C: C: \textbackslash{}( \textbackslash{}sin\textasciicircum{}\{-1\}\textbackslash{}left[n\textbackslash{}sin\textbackslash{}left(\textbackslash{}sin\textasciicircum{}\{-1\}\textbackslash{}frac\{L\}\{R\} - \textbackslash{}sin\textasciicircum{}\{-1\}\textbackslash{}frac\{L\}\{nR\}\textbackslash{}right)\textbackslash{}right] \textbackslash{})

\paragraph{Figure/image path.}
/root/proposal\_for\_physic/science-mango/phyx\_mini\_run/phyx\_data/test\_image/18.png

\paragraph{Formalization target.}
create a compiling Lean file with sorry bodies at `PhyXMiniProblems/problem\_phyx\_mini\_0018.lean`. The Lean declarations must preserve the physical quantities, dimensions or dimensional roles, figure labels, governing-law hypotheses, and final relation expressed by this problem.
Use Mathlib/Physlib names found through LeanExplore where available. If a physics API is missing, introduce faithful local abstractions rather than scalar placeholder aliases.

\begin{theorem}[Physics formalization target]
\label{thm:physics:phyx_mini_0018:target}
\lean{PhyXMiniProblems.ProblemPhyXMini0018.problem_phyx_mini_0018}
\uses{def:physics:phyx-mini-0018:phyxminiproblems-problemphyxmini0018-lengthinmeters, def:physics:phyx-mini-0018:phyxminiproblems-problemphyxmini0018-quartercirclerefractionsetup, def:physics:phyx-mini-0018:phyxminiproblems-problemphyxmini0018-hasphysicalparameters, def:physics:phyx-mini-0018:phyxminiproblems-problemphyxmini0018-hasphysicalrayangles, def:physics:phyx-mini-0018:phyxminiproblems-problemphyxmini0018-obeysquartercirclegeometry, def:physics:phyx-mini-0018:phyxminiproblems-problemphyxmini0018-satisfiessnelllawat}
For a horizontal ray entering the quarter-circular material at height \(L\),
the acute emergence angle is
\[
\arcsin\!\left(
 n\sin\!\left(\arcsin\frac LR-\arcsin\frac{L}{nR}\right)\right),
\]
which is answer C.
\end{theorem}
\begin{proof}
The radius geometry gives the curved-face incidence angle
\(\alpha=\arcsin(L/R)\).  Snell's law from vacuum into the material gives
the internal normal angle \(\beta=\arcsin(L/(nR))\).  The ray therefore meets
the flat face at angle \(\alpha-\beta\); applying Snell's law there and
selecting the acute emergence branch yields the displayed formula.
\end{proof}

% --- Archon declaration topology begin ---
\section{Declaration topology}
\begin{definition}[Dim Length]
  \label{def:physics:phyx-mini-0018:phyxminiproblems-problemphyxmini0018-dimlength}
  \lean{PhyXMiniProblems.ProblemPhyXMini0018.DimLength}
  A physical length, independent of the unit used to read its value
\end{definition}
\begin{proof}
  This is the declared model definition of Dim Length.
\end{proof}
\begin{definition}[length In Meters]
  \label{def:physics:phyx-mini-0018:phyxminiproblems-problemphyxmini0018-lengthinmeters}
  \lean{PhyXMiniProblems.ProblemPhyXMini0018.lengthInMeters}
  \uses{def:physics:phyx-mini-0018:phyxminiproblems-problemphyxmini0018-dimlength}
  The scalar meter readout of a dimensionful length
\end{definition}
\begin{proof}
  This is the declared model definition of length In Meters.
\end{proof}
\begin{definition}[Optical Medium]
  \label{def:physics:phyx-mini-0018:phyxminiproblems-problemphyxmini0018-opticalmedium}
  \lean{PhyXMiniProblems.ProblemPhyXMini0018.OpticalMedium}
  A homogeneous optical medium together with its positive, dimensionless refractive index
\end{definition}
\begin{proof}
  This is the declared model definition of Optical Medium.
\end{proof}
\begin{definition}[Optical Interface]
  \label{def:physics:phyx-mini-0018:phyxminiproblems-problemphyxmini0018-opticalinterface}
  \lean{PhyXMiniProblems.ProblemPhyXMini0018.OpticalInterface}
  The two material interfaces crossed by the ray in the figure
\end{definition}
\begin{proof}
  This is the declared model definition of Optical Interface.
\end{proof}
\begin{definition}[Quarter Circle Refraction Setup]
  \label{def:physics:phyx-mini-0018:phyxminiproblems-problemphyxmini0018-quartercirclerefractionsetup}
  \lean{PhyXMiniProblems.ProblemPhyXMini0018.QuarterCircleRefractionSetup}
  \uses{def:physics:phyx-mini-0018:phyxminiproblems-problemphyxmini0018-dimlength, def:physics:phyx-mini-0018:phyxminiproblems-problemphyxmini0018-opticalmedium}
  Physical quantities and angle labels for the quarter-circle ray diagram. The body is the quarter disk of radius radius; incomingHeightAboveBase is the figure label L. The internal and outgoing directions are recorded as positive angles below the horizontal base direction, matching the dashed horizontal normal at the flat exit face
\end{definition}
\begin{proof}
  This is the declared model definition of Quarter Circle Refraction Setup.
\end{proof}
\begin{definition}[Has Physical Parameters]
  \label{def:physics:phyx-mini-0018:phyxminiproblems-problemphyxmini0018-hasphysicalparameters}
  \lean{PhyXMiniProblems.ProblemPhyXMini0018.HasPhysicalParameters}
  \uses{def:physics:phyx-mini-0018:phyxminiproblems-problemphyxmini0018-lengthinmeters, def:physics:phyx-mini-0018:phyxminiproblems-problemphyxmini0018-quartercirclerefractionsetup}
  The nondegenerate material and length conditions in the pictured situation. The surrounding vacuum has index one and the material is optically denser
\end{definition}
\begin{proof}
  This is the declared model definition of Has Physical Parameters.
\end{proof}
\begin{definition}[Has Physical Ray Angles]
  \label{def:physics:phyx-mini-0018:phyxminiproblems-problemphyxmini0018-hasphysicalrayangles}
  \lean{PhyXMiniProblems.ProblemPhyXMini0018.HasPhysicalRayAngles}
  \uses{def:physics:phyx-mini-0018:phyxminiproblems-problemphyxmini0018-quartercirclerefractionsetup}
  All four normal-based ray angles use their acute physical branches. These conditions select the intended inverse-sine branches and assert that the ray emerges rather than undergoing total internal reflection
\end{definition}
\begin{proof}
  This is the declared model definition of Has Physical Ray Angles.
\end{proof}
\begin{definition}[Obeys Quarter Circle Geometry]
  \label{def:physics:phyx-mini-0018:phyxminiproblems-problemphyxmini0018-obeysquartercirclegeometry}
  \lean{PhyXMiniProblems.ProblemPhyXMini0018.ObeysQuarterCircleGeometry}
  \uses{def:physics:phyx-mini-0018:phyxminiproblems-problemphyxmini0018-lengthinmeters, def:physics:phyx-mini-0018:phyxminiproblems-problemphyxmini0018-quartercirclerefractionsetup}
  Geometry read from the quarter-circle figure. The incoming ray is parallel to the base. The radius through its entry point has vertical component L, so the entry incidence angle is arcsin (L / R). Subtracting the inside-normal angle gives the ray's angle below the horizontal flat-face normal
\end{definition}
\begin{proof}
  This is the declared model definition of Obeys Quarter Circle Geometry.
\end{proof}
\begin{definition}[Satisfies Snell Law At]
  \label{def:physics:phyx-mini-0018:phyxminiproblems-problemphyxmini0018-satisfiessnelllawat}
  \lean{PhyXMiniProblems.ProblemPhyXMini0018.SatisfiesSnellLawAt}
  \uses{def:physics:phyx-mini-0018:phyxminiproblems-problemphyxmini0018-opticalinterface, def:physics:phyx-mini-0018:phyxminiproblems-problemphyxmini0018-quartercirclerefractionsetup}
  Snell's law at either interface: refractive index times the sine of the angle from the interface normal is conserved
\end{definition}
\begin{proof}
  This is the declared model definition of Satisfies Snell Law At.
\end{proof}
\begin{lemma}[curved Refraction Angle eq arcsin]
  \label{lem:physics:phyx-mini-0018:phyxminiproblems-problemphyxmini0018-curvedrefractionangle-eq-arcsin}
  \lean{PhyXMiniProblems.ProblemPhyXMini0018.curvedRefractionAngle_eq_arcsin}
  \uses{def:physics:phyx-mini-0018:phyxminiproblems-problemphyxmini0018-lengthinmeters, def:physics:phyx-mini-0018:phyxminiproblems-problemphyxmini0018-quartercirclerefractionsetup, def:physics:phyx-mini-0018:phyxminiproblems-problemphyxmini0018-hasphysicalparameters, def:physics:phyx-mini-0018:phyxminiproblems-problemphyxmini0018-hasphysicalrayangles, def:physics:phyx-mini-0018:phyxminiproblems-problemphyxmini0018-obeysquartercirclegeometry, def:physics:phyx-mini-0018:phyxminiproblems-problemphyxmini0018-satisfiessnelllawat}
  The entry-face geometry and Snell's law determine the in-material angle β = arcsin (L / (n R)). This is a derived intermediate result, not a governing-law assumption
\end{lemma}
\begin{proof}
  The conclusion follows from the governing relations and figure readouts named in the statement of curved Refraction Angle eq arcsin.
\end{proof}
% --- Archon declaration topology end ---
% --- Archon physics formalization source end ---
